\section{Physics-first rewrite patch set}

\paragraph{Scope.}
This note is a patch set for rewriting the manuscript in physical-first
order without editing \texttt{proj.tex}.  It uses the current anchors in
\texttt{proj.tex} and the prior rewrite map in
\texttt{agent\_material/08\_landmark\_physics\_rewrite\_plan.tex}.  The
owned output is this file only.  Each block below is written so that the
claim strength is either mathematically stronger than the current text or
epistemically sharper.

\subsection{Claim attacked}

The remaining global failure mode is not the Borcherds product itself.
The product, the Gritsenko--Nikulin denominator theorem, and the
Oberdieck--Pixton scalar square are already quarantined with usable
status.  The failure mode is the order of explanation: several passages
still let a lattice or modular object appear before the physical
charge/index operation which justifies using it.  The replacement policy is:
\[
\hbox{physical protected index}
\longrightarrow
\hbox{Gram-invariant index lattice}
\longrightarrow
\hbox{automorphic product and Weyl chamber}
\longrightarrow
\hbox{denominator algebra}
\longrightarrow
\hbox{formal current/Hall problem}.
\]
At no point may \(\Gamma_{\mathrm{gram}}\cong \mathbb Z^3\) be called the
full microscopic BPS charge lattice unless a compactification charge
lattice and projection theorem have been supplied.  At no point may a
denominator identity be called a physical BPS bracket.

\subsection{Patch index}

\begin{center}
\begin{tabular}{c|p{0.32\linewidth}|p{0.43\linewidth}}
Patch & Anchor in \texttt{proj.tex} & Upgrade \\
\hline
P1 & 55--115, abstract and opening & Replace modular-first summary by
protected-index-first summary.  Statuses: OP scalar square proved; virtual
determinant constructed; BKM denominator proved as denominator algebra;
physical BPS bracket open.\\
P2 & Body order, especially 120--211, 213--404, 542--709 & Move the OP
scalar theorem and K3 half-index construction before the duality and Weyl
sections.  The automorphic formalism becomes a consequence of the
Gram-reduced index datum, not an input.\\
P3 & 120--211, charge lattice and duality & Replace ``charge lattice'' by
``Gram-invariant index lattice'' and make the Harvey--Moore/DVV reduction an
explicit physical assumption.\\
P4 & 215--220, one-particle product lattice & Replace
\(\Gamma_{\mathrm{BPS}}\) by \(\Gamma_{\mathrm{gram}}\), with no microscopic
bracket claim.\\
P5 & 840--862, exterior-square duality datum & Replace ``geometric duality
group'' by ``automorphic duality datum''.  The group isomorphism is a
mathematical theorem; its physical-duality interpretation is an input.\\
P6 & 866--870 and 1031--1054, Weyl chamber and charge degrees & Replace
``same BPS charge lattice'' and ``BPS charge grading'' by Gram-degree
language.  Active support is an index-support condition.\\
P7 & 1211--1475, denominator/current/Hall transition & Rename the theorem
status from bracket-level to denominator-level.  State explicitly that Hall
or factorization data are required before there is a physical BPS algebra.\\
P8 & 1681--1693, determinant forgetfulness & Replace ``one-particle BPS
index coefficient'' by ``virtual K3 half-index coefficient / signed root
superdimension''.\\
P9 & 1893--3406, late rows and Vol.~III comparisons & Keep these sections
late.  They are formal product/current or imported denominator rows, not
input to the \(N=1\) physical derivation and not physical row maps.\\
\end{tabular}
\end{center}

\subsection{P1. Abstract replacement}

\paragraph{Anchor.}
Replace \texttt{proj.tex} lines 55--115.

\paragraph{Replacement block.}
\begin{verbatim}
\begin{abstract}
Let \(X=S\times E\), with \(S\) a K3 surface and \(E\) an elliptic
curve.  The first input in this paper is physical and enumerative, not an
a priori choice of a Siegel modular form: a protected two-charge index in a
cusp polarization, reduced to the three Gram invariants
\[
        n={1\over 2}(Q,Q),\qquad l=(Q,P),\qquad
        m={1\over 2}(P,P).
\]
The reduction from the full compactification charge lattice to these
three invariants is used as a Harvey--Moore/Dijkgraaf--Verlinde--Verlinde
index assumption.  The resulting lattice
\(\Gamma_{\mathrm{gram}}\cong \mathbb Z^3\) is the universal
Gram-invariant index lattice of the product.  It is not asserted to be the
full microscopic BPS charge lattice.

Oberdieck--Pixton prove the reduced primitive curve-counting identity
\[
        Z_{\mathrm{OP}}(q,r,s)=-\chi_{10}(Z)^{-1}
        =-4096\,\Delta_5(Z)^{-2}.
\]
Thus the enumerative scalar partition function is the square of the
normalized Igusa square root \(\Delta_5\), up to the conventional factor
\(-4096\).  The square-root object used below is algebraic and virtual:
it is the Borcherds determinant built from the virtual Ramond-sector K3
half-index character
\[
        \operatorname{sch}(V_{\mathrm{vir}})=\phi_{0,1}=Z_{\mathrm{K3}}/2
\]
in the Grothendieck group of \(\mathbb Z^2\)-graded super vector spaces.
With \(c(D)\) denoting the Fourier coefficients of \(\phi_{0,1}\), the
determinant is
\[
D_X(q,r,s)=q^{1/2}r^{1/2}s^{1/2}
\prod_{\substack{n,m\geq 0,\ l\in\mathbb Z\\ (n,l,m)>0}}
        (1-q^n r^l s^m)^{c(4nm-l^2)}
        =\Delta_5(Z).
\]

The Borcherds--Kac--Moody superalgebra used in the paper is proved only
at denominator level.  Its signed root superdimensions are the same
coefficients \(c(4nm-l^2)\), and its Weyl--Kac--Borcherds denominator
identity is the automorphic product above.  The Harvey--Moore
interpretation identifies the grading of this denominator identity with
the protected index grading after the Gram-reduction assumption.  It does
not by itself construct a Lie bracket on microscopic BPS states.  The
bracket-level BPS algebra, and any Hall or factorization algebra
realizing it, remain separate open constructions unless explicitly
supplied.
\end{abstract}
\end{verbatim}

\paragraph{Status label.}
Proved: OP scalar identity and Borcherds determinant identity in the cited
normalization.  Constructed: the virtual half-index determinant.  Assumed:
the physical Gram-reduction of the protected index.  Open: microscopic
BPS bracket and Hall/factorization realization.

\subsection{P2. Replacement theorem order}

\paragraph{Anchor.}
Use this order for the next manuscript rewrite.  The line numbers refer to
the current \texttt{proj.tex}.

\paragraph{Replacement order.}
\begin{enumerate}
\item Opening and abstract: P1.
\item Protected enumerative scalar theorem: move the material now at
lines 542--709 immediately after the introduction.  This places the
Oberdieck--Pixton theorem before any lattice, Weyl, or modular-form
classification.
\item K3 Ramond-sector and virtual half-index: keep the material now at
lines 222--404 after the scalar theorem.  It explains why the square-root
determinant is virtual and why no honest Hilbert-space square root has
been constructed.
\item Charge-to-index reduction: place the replacement P3 after the
half-index.  The physical input is now the charge/index operation, not the
Siegel form.
\item Automorphic square root and exterior-square linear algebra: place
the material now at lines 740--862 after the charge/index reduction, with
P5 replacing the overclaimed duality language.
\item Weyl chamber and Borcherds cone: place the material now at
lines 864--1125 after the automorphic datum, with P6 replacing BPS-lattice
language by Gram-degree language.
\item Denominator algebra: place the material now at lines 1211--1321
after the Weyl chamber.  Use P7 to mark it as denominator algebra, not
physical state algebra.
\item Hall/factorization realization problem: keep the material now at
lines 1373--1475 after the denominator algebra, again with P7.  This is
the first location where bracket-level physical BPS algebra may be named.
\item Determinant-forgetfulness and late extensions: keep the material now
at lines 1658--3406 late, using P8 and P9 to prevent late comparison rows
from functioning as hidden inputs.
\end{enumerate}

\subsection{P3. Charge lattice and duality replacement}

\paragraph{Anchor.}
Replace \texttt{proj.tex} lines 120--211, including the current charge
matrix convention and BPS charge pairing definitions.

\paragraph{Replacement block.}
\begin{verbatim}
\section{Protected charges and the Gram-invariant index lattice}

The physical object used in the paper is a protected index for a primitive
electric--magnetic charge pair.  The product formula will use only the
three Gram invariants of that pair.  This section records the precise
point at which physical input enters.

\begin{assumption}[Protected two-charge index and Gram reduction]
\label{ass:protected-gram-reduction}
Fix a cusp polarization of the compactification and a primitive pair of
charges \((Q,P)\) in an even integral charge lattice
\(\Lambda_{\mathrm{phys}}\).  Let
\[
        n={1\over2}(Q,Q),\qquad l=(Q,P),\qquad
        m={1\over2}(P,P).
\]
In the chamber under consideration, the protected index
\(\Omega(Q,P)\) depends only on \((n,l,m)\).  Equivalently, there is a
function \(\Omega_{\mathrm{gram}}\) such that
\[
        \Omega(Q,P)=\Omega_{\mathrm{gram}}(n,l,m).
\]
The electric--magnetic \(SL_2(\mathbb Z)\)-action on \((Q,P)\) acts on
the Gram matrix
\[
        G(Q,P)=\begin{pmatrix}2n&l\\ l&2m\end{pmatrix}
\]
by \(G\mapsto AGA^t\).  The enlargement from this \(SL_2\)-covariance to
the Siegel modular group used by the automorphic product is the standard
Harvey--Moore/Dijkgraaf--Verlinde--Verlinde duality input, not a theorem
proved from the compactification in this paper.
\end{assumption}

\begin{definition}[Gram-invariant index lattice]
\label{def:gram-index-lattice}
The lattice used by the product is
\[
        \Gamma_{\mathrm{gram}}=\mathbb Z^3
        =\{(n,l,m): n,l,m\in\mathbb Z\}.
\]
It is the lattice of Gram invariants of the protected index.  Its
discriminant is
\[
        D(n,l,m)=4nm-l^2.
\]
The cusp-positive semigroup is
\[
        (n,l,m)>0
        \quad\Longleftrightarrow\quad
        m>0,\quad\hbox{or }m=0,n>0,\quad\hbox{or }m=n=0,l<0.
\]
This is an index semigroup.  It is not the full microscopic charge
lattice unless an additional theorem identifies the compactification
charges with this quotient.
\end{definition}

\begin{proposition}[Discriminant and character convention]
\label{prop:gram-discriminant-character}
Under Assumption~\ref{ass:protected-gram-reduction}, the variables in the
Fourier expansion of the protected index are indexed by
\(\Gamma_{\mathrm{gram}}\):
\[
        q^n r^l s^m
        =\exp(2\pi i(n\tau+lz+m\sigma)).
\]
The sign and parity data used by the Borcherds product depend only on the
discriminant \(D=4nm-l^2\) and on the multiplier character of the
automorphic form.  This statement is bookkeeping after the
Gram-reduction assumption; it is not a construction of a BPS bracket.
\end{proposition}

\begin{remark}[Physical status of the duality group]
\label{rem:duality-status}
The arithmetic isomorphism
\(\PSp_4(\mathbb Z)\simeq \SO^+(\Lambda^{3,2})\) is a mathematical
statement about the automorphic lattice.  Its interpretation as the
physical duality group of the compactification is used only with the
status of the Harvey--Moore/DVV physical input above.
\end{remark}
\end{verbatim}

\paragraph{Status label.}
Assumption: protected index depends only on the Gram invariants and has
the stated duality covariance.  Proved after the assumption: the
discriminant bookkeeping and the \(\Gamma_{\mathrm{gram}}\) product
indexing.  Refused: identifying \(\Gamma_{\mathrm{gram}}\) with the full
microscopic charge lattice.

\subsection{P4. One-particle determinant preamble replacement}

\paragraph{Anchor.}
Replace \texttt{proj.tex} lines 215--220.

\paragraph{Replacement block.}
\begin{verbatim}
The determinant is formed over the Gram-invariant index lattice
\[
        \Gamma_{\mathrm{gram}}=\mathbb Z^3,\qquad
        \gamma=(n,l,m).
\]
The chamber order is
\[
\gamma>0\quad\Longleftrightarrow\quad
m>0,\quad\hbox{or }m=0,n>0,\quad\hbox{or }m=n=0,l<0.
\]
This is the positive cone for the protected-index product.  It is not an
assertion that \(\mathbb Z^3\) is the full BPS charge lattice of the
compactification.
\end{verbatim}

\paragraph{Status label.}
Constructed: determinant indexing.  Refused: microscopic charge-lattice
identification.

\subsection{P5. Exterior-square duality replacement}

\paragraph{Anchor.}
Replace \texttt{proj.tex} lines 840--862.

\paragraph{Replacement block.}
\begin{verbatim}
\begin{definition}[Automorphic duality datum]
\label{def:automorphic-duality-datum}
The automorphic duality datum used in the paper is the arithmetic action
of
\[
        \PSp_4(\mathbb Z)
\]
on the Siegel upper half-space, transported by the exterior-square
construction to the index-two orthogonal group
\[
        \SO_+(\Lambda^{3,2})\subset \Orth(\Lambda^{3,2}).
\]
Equivalently, after quotienting by the central sign,
\[
        \PSp_4(\mathbb Z)\simeq \SO_+(\Lambda^{3,2}).
\]
This is a theorem of the exterior-square construction and integral
lattice arithmetic.  The statement that this automorphic group is the
physical electric--magnetic duality group of the compactification is not
proved here; it is part of the Harvey--Moore/DVV duality input attached
to the protected index.
\end{definition}
\end{verbatim}

\paragraph{Status label.}
Proved: exterior-square arithmetic.  Assumed: physical-duality
interpretation.

\subsection{P6. Weyl chamber and Gram-degree replacements}

\paragraph{Anchor A.}
Replace \texttt{proj.tex} lines 866--870.

\paragraph{Replacement block A.}
\begin{verbatim}
We now pass from the Gram-invariant index lattice to the Lorentzian
lattice used by the Borcherds denominator.  The cusp polarization fixes a
positive semigroup in \(\Gamma_{\mathrm{gram}}\); the corresponding Weyl
chamber is a chamber for the automorphic product.  A microscopic
wall-crossing interpretation requires additional physical input and is not
proved in this section.
\end{verbatim}

\paragraph{Anchor B.}
Replace \texttt{proj.tex} lines 1031--1054.

\paragraph{Replacement block B.}
\begin{verbatim}
\begin{definition}[Gram degrees in the Weyl chamber]
\label{def:gram-degrees-weyl}
For \(\gamma=(n,l,m)\in\Gamma_{\mathrm{gram}}\), set
\[
        \alpha(\gamma)=(2n,-l,2m)\in\Lambda^{2,1}_{II}.
\]
The product monomial \(q^n r^l s^m\) is the character of this
Gram-degree.  The degree is active if
\[
        c(4nm-l^2)\neq 0.
\]
Thus activity is a statement about the support of the protected index, or
equivalently about signed root superdimension in the denominator algebra.
It is not, by itself, a statement that a microscopic BPS state or BPS
bracket has been constructed in that degree.
\end{definition}
\end{verbatim}

\paragraph{Status label.}
Proved: Weyl chamber arithmetic and active-support computation.  Refused:
physical wall-crossing or microscopic state-algebra interpretation without
a construction.

\subsection{P7. Denominator/current/Hall transition replacement}

\paragraph{Anchor A.}
Replace the opening of the denominator section at
\texttt{proj.tex} lines 1211--1223.

\paragraph{Replacement block A.}
\begin{verbatim}
\section{The denominator algebra}

The previous sections identify \(\Delta_5\) as a virtual determinant and
as a Borcherds automorphic product.  We now use the associated
Borcherds--Kac--Moody superalgebra only with denominator-level status.
A denominator identity specifies real roots, imaginary simple roots,
signed root superdimensions, and the Weyl anti-invariant product.  It
does not, by itself, construct the bracket on microscopic BPS states.

For this reason the algebra below is called the denominator algebra.  A
physical BPS algebra is a stronger object: it requires a construction of
states, a charge grading, and a bracket or Hall product whose signed
character is the denominator.
\end{verbatim}

\paragraph{Anchor B.}
Replace \texttt{proj.tex} lines 1254--1278.  If cross-reference churn is
undesirable during integration, the old label may be kept as a temporary
alias, but the theorem name and text must remain denominator-level.

\paragraph{Replacement block B.}
\begin{verbatim}
\begin{theorem}[Denominator algebra square root]
\label{thm:denominator-algebra-square-root}
The denominator Lie superalgebra has the following Gram-degree
decomposition:
\[
\mathcal A_{\mathrm{den}}
=\bigoplus_{\gamma\in\Gamma_{\mathrm{gram}}}
\mathcal A_{\mathrm{den},\gamma},
\qquad
[\mathcal A_{\mathrm{den},\gamma},\mathcal A_{\mathrm{den},\gamma'}]
\subset \mathcal A_{\mathrm{den},\gamma+\gamma'}.
\]
Its signed Gram-degree multiplicities are the K3 virtual half-index
coefficients:
\[
\sdim\mathcal A_{\mathrm{den},(n,l,m)}
=f(nm,l).
\]
Its Weyl--Kac--Borcherds denominator is
\[
\operatorname{den}(\mathcal A_{\mathrm{den}})
=64^{-1}\Delta_5(2Z).
\]
This is a theorem about the Borcherds denominator algebra.  It is not a
construction of the microscopic BPS Hilbert space, BPS bracket, or
operator product.
\end{theorem}
\end{verbatim}

\paragraph{Anchor C.}
Replace the Hall criterion condition beginning at
\texttt{proj.tex} lines 1373--1384.

\paragraph{Replacement block C.}
\begin{verbatim}
\begin{theorem}[Microscopic Hall realization criterion]
\label{thm:microscopic-hall-realization-criterion}
A bracket-level BPS algebra would require the following additional data.
\begin{enumerate}
\item A compactification charge lattice \(\Lambda_{\mathrm{phys}}\), a
protected index on primitive charge pairs, and a projection to
\(\Gamma_{\mathrm{gram}}\) satisfying
Assumption~\ref{ass:protected-gram-reduction}.
\item A moduli problem or factorization category whose objects carry this
charge grading and whose Hall product is defined in the chamber under
consideration.
\item A proof that the signed graded character of the Hall or
factorization algebra is the coefficient system
\(c(4nm-l^2)\).
\item A proof that the resulting bracket satisfies the Borcherds
relations for the same real roots and imaginary simple roots as the
denominator algebra.
\end{enumerate}
Until these data are supplied, the manuscript has a denominator algebra
and a formal current envelope, not a physical BPS bracket.
\end{theorem}
\end{verbatim}

\paragraph{Anchor D.}
Replace \texttt{proj.tex} lines 1446--1456.

\paragraph{Replacement block D.}
\begin{verbatim}
\begin{openproblem}[Microscopic square root of the primitive scalar]
\label{op:microscopic-square-root-primitive-scalar}
Oberdieck--Pixton fix the primitive reduced curve-counting scalar
\[
        Z_{\mathrm{OP}}=-4096\,\Delta_5^{-2}.
\]
After the Harvey--Moore/DVV dictionary, this scalar is expected to be the
square of the protected-index character in the primitive chamber.  The
missing theorem is stronger: construct a BPS or Hall/factorization object
whose signed character is \(\Delta_5^{-1}\), whose denominator algebra is
the \(\Delta_5\) Borcherds algebra, and whose completed square recovers
the OP scalar.  The OP theorem alone does not construct that object.
\end{openproblem}
\end{verbatim}

\paragraph{Status label.}
Proved: denominator identity and formal current character.  Formal:
current envelope.  Open: physical Hall/factorization algebra and bracket.
Refused: saying OP alone fixes the protected character.

\subsection{P8. Determinant-forgetfulness replacement}

\paragraph{Anchor.}
Replace \texttt{proj.tex} lines 1681--1693.

\paragraph{Replacement block.}
\begin{verbatim}
\begin{proposition}[What the determinant forgets]
\label{prop:determinant-forgets-bracket}
The exponent \(c(4nm-l^2)\) in the determinant is the signed coefficient
of the virtual K3 half-index.  In the denominator algebra it is the signed
root superdimension of the corresponding root degree.  The determinant
therefore remembers the graded Euler character and the automorphic
denominator identity.

It forgets the stronger microscopic data: a basis of BPS states, a
chain-level moduli problem, and a bracket or Hall product realizing the
same coefficients.  Recovering those data is exactly the bracket-level
BPS algebra problem.
\end{proposition}
\end{verbatim}

\paragraph{Status label.}
Proved: determinant remembers signed graded dimensions.  Open: recovery
of microscopic state and bracket data.

\subsection{P9. Late extension quarantine}

\paragraph{Anchor A.}
Insert at the beginning of the late extension material, before the
section now starting at \texttt{proj.tex} line 1893.

\paragraph{Insertion block A.}
\begin{verbatim}
\paragraph{Status of the extension sections.}
The following sections are extension and comparison material.  They are
not inputs to the \(N=1\) derivation of \(\Delta_5\) from the protected
K3 index and the OP scalar theorem.  For each row, the automorphic product
or denominator identity has exactly the status stated locally.  A physical
row map or BPS bracket exists only when a compactification sector,
protected-index map, and Hall/factorization construction have been
supplied.
\end{verbatim}

\paragraph{Anchor B.}
Insert before the Vol.~III row-map comparison now around
\texttt{proj.tex} lines 3221--3406.

\paragraph{Insertion block B.}
\begin{verbatim}
\paragraph{Vol.~III row-map status.}
The symbols \(F_i\) from the Clery--Gritsenko catalogue and the physical
Vol.~III symbols \(\Phi^{\mathrm{V3}}_j\) belong to different row systems
until a row map has been constructed.  The comparison below may record
weights, characters, multiplier covers, and imported CHL/WKB denominator
theorems.  It must not be used as a physical derivation of any row unless
the map fixes the compactification sector, the protected index, the
charge projection, and the bracket or Hall realization.
\end{verbatim}

\paragraph{Status label.}
Formal/product status for comparison rows unless a local theorem upgrades
them.  Imported denominator status for Govindarajan CHL/WKB rows as
stated in the manuscript.  Open: physical row maps and boundary-row
bracket constructions.

\subsection{Broken statements to upgrade}

\begin{enumerate}
\item \texttt{proj.tex:170}: ``The charge lattice is even'' should become
``The compactification charge lattice is assumed even at the point of
physical input; the product uses only the induced Gram-invariant lattice.''
\item \texttt{proj.tex:215}: ``The charge lattice is
\(\Gamma_{\mathrm{BPS}}=\mathbb Z^3\)'' should become P4.
\item \texttt{proj.tex:841}: ``The geometric duality group is
\(\PSp_4(\mathbb Z)\)'' should become P5.
\item \texttt{proj.tex:866}: ``same BPS charge lattice'' should become
``same Gram-invariant index lattice'' or the fuller P6.
\item \texttt{proj.tex:1381}: ``The protected charge lattice is
\(\Gamma_{\mathrm{BPS}}\)'' should become the projected-index condition in
P7.
\item \texttt{proj.tex:1450}: ``OP fixes the protected character'' should
be replaced by P7D.
\item \texttt{proj.tex:1681--1693}: ``one-particle BPS index
coefficient'' should become P8.
\item \texttt{proj.tex:2008, 3253, 3272, 3283, 3404}: physical row-map
language should remain only as explicitly open row-map status, as in P9.
\end{enumerate}

\subsection{Claims strengthened}

\begin{enumerate}
\item The OP theorem is isolated as a proved scalar identity for primitive
reduced curve counting, not as a direct construction of BPS states.
\item The square root is identified as the normalized Borcherds
determinant of a virtual K3 half-index, so the square-root claim no longer
depends on an unbuilt Hilbert-space square root.
\item \(\Gamma_{\mathrm{gram}}\cong\mathbb Z^3\) is given the exact status
of a Gram-invariant index lattice.
\item \(\PSp_4(\mathbb Z)\simeq\SO_+(\Lambda^{3,2})\) is kept as a
mathematical automorphic theorem, while its compactification-duality role
is explicitly an input.
\item The denominator algebra is separated from bracket-level physical
BPS algebra.
\item Late row comparisons are prevented from serving as hidden
physical-origin arguments.
\end{enumerate}

\subsection{Claims refused}

\begin{enumerate}
\item Refused: \(\Gamma_{\mathrm{BPS}}=\mathbb Z^3\) as the microscopic
BPS charge lattice.
\item Refused: \(\PSp_4(\mathbb Z)\) as a physical duality theorem proved
by exterior-square linear algebra alone.
\item Refused: OP scalar square as a construction of the protected BPS
character before the Harvey--Moore/DVV dictionary.
\item Refused: denominator identity as physical BPS bracket.
\item Refused: active Borcherds support as a microscopic state-support
theorem.
\item Refused: Clery--Gritsenko or Vol.~III row labels as physical row
maps without a constructed map.
\end{enumerate}

\subsection{Formulas checked}

The local scripts checked the formulas needed by this patch set:
\begin{enumerate}
\item \texttt{compute/verify\_lattice.py}: the Pfaffian Gram matrix in
the \((f_1,f_2,f_3,f_{-2},f_{-1})\) basis; the type-II simple-root Gram
matrix with diagonal \(2\) and off-diagonal \(-2\); the Weyl vector
\(\rho=(1,-1/2,1)\); the pairings
\((z_1,z_2,z_3)=(-2n,-2l,-2m)\); and the half-root relation used in the
denominator normalization.
\item \texttt{compute/verify\_square\_root.py}: the coefficients of
\(\phi_{0,1}\) through \(q^2\),
\[
q^0:\ y^{-1}+10+y,\qquad
q^1:\ 10y^{-2}-64y^{-1}+108-64y+10y^2,
\]
and the first coefficients of \(\prod_{k\ge1}(1-q^k)^{9}\), namely
\[
1,-9,27,-12,-90,135,54,-99,-189,-85.
\]
\end{enumerate}

\subsection{Commands run}

\begin{verbatim}
sed -n '1,220p' AGENTS.md
sed -n '1,200p' CLAUDE.md
sed -n '128,148p' /Users/raeez/ecosystem/INVARIANTS.md
sed -n '133,148p' /Users/raeez/ecosystem/AGENTS-HARNESS.md
sed -n '1,430p' agent_material/08_landmark_physics_rewrite_plan.tex
nl -ba proj.tex | sed -n '50,225p'
nl -ba proj.tex | sed -n '226,410p'
nl -ba proj.tex | sed -n '542,736p'
nl -ba proj.tex | sed -n '739,1125p'
nl -ba proj.tex | sed -n '1210,1475p'
nl -ba proj.tex | sed -n '1658,1735p'
sed -n '1,240p' compute/verify_lattice.py
sed -n '1,260p' compute/verify_square_root.py
sed -n '1,220p' /Users/raeez/calabi-yau-quantum-groups/FRONTIER.md
python3 compute/verify_lattice.py
python3 compute/verify_square_root.py
rg -n -F -e "charge lattice is" -e "same BPS charge lattice" \
  -e "geometric duality group" -e "protected character" \
  -e "one-particle BPS index coefficient" -e "BPS charge lattice" \
  -e "physical row" -e "Gamma_\\mathrm{BPS}" \
  -e "BPS charge grading" proj.tex agent_material
rg -n -e "Conjectural|CHL|Vol III|Gritsenko--Clery|..." proj.tex
rg -n -e "Borcherds|Gritsenko|Nikulin|Oberdieck|..." proj.bib
rg -n "openproblem|assumption|newtheorem|theoremstyle" \
  raeez-math-template.sty proj.tex
perl -ne 'print "$.:$_" if /[^\x00-\x7F]/' \
  agent_material/09_physics_first_rewrite_patchset.tex
chktex -q agent_material/09_physics_first_rewrite_patchset.tex || true
git status --short
git status --short -- agent_material/09_physics_first_rewrite_patchset.tex
\end{verbatim}

\subsection{Remaining open obligations}

\begin{enumerate}
\item Prove or cite a compactification theorem constructing the full
physical charge lattice \(\Lambda_{\mathrm{phys}}\) and the protected
index projection to \(\Gamma_{\mathrm{gram}}\).
\item Prove the Harvey--Moore/DVV duality dictionary at the exact
normalization used here, or keep it as an explicit assumption.
\item Construct an honest Hilbert-space or derived-category square root
whose signed character is \(\Delta_5^{-1}\), not only the virtual
half-index determinant.
\item Construct the Hall/factorization object whose bracket realizes the
\(\Delta_5\) denominator algebra.
\item Prove the physical row maps for the late Clery--Gritsenko and
Vol.~III comparison rows before using them as physical derivations.
\end{enumerate}
